\documentclass[12pt,a4paper]{article}
\usepackage[utf8]{inputenc}
\usepackage[T5]{fontenc}
\usepackage[vietnamese]{babel}
\usepackage{graphicx}
\usepackage{geometry}
\geometry{a4paper, margin=2cm}
\usepackage{tikz}
\usetikzlibrary{shapes.geometric, arrows, positioning}
\usepackage{hyperref}
\usepackage{enumitem}

\title{\textbf{BÁO CÁO NHIỆM VỤ TUẦN 1}\\ Phát triển AI Agent phát hiện bệnh trên lá lúa}
\author{}
\date{}

\begin{document}

\maketitle

\section{Task 1: Review các bài báo khoa học và Mô hình đề tài sử dụng}

\subsection{1. Review các bài báo khoa học có cùng chủ đề}

Dưới đây là phần tổng hợp và đánh giá chi tiết 3 bài báo tiêu biểu liên quan đến nhận diện và phân loại bệnh trên lá lúa:

\textbf{Bài báo 1: Neural Network-based Study for Rice Leaf Disease Recognition and Classification: A Comparative Analysis Between Feature-based Model and Direct Imaging Model}
\begin{itemize}
    \item \textbf{Mô hình sử dụng:} Mô hình mạng nơ-ron nhân tạo (ANN), cụ thể phân tích so sánh giữa Mô hình phát hiện phân tích đặc trưng (FADM) với các thuật toán trích xuất đặc trưng (GLCM, DWT...) kết hợp giảm chiều dữ liệu (PCA, KPCA) và Mô hình nhận diện trực tiếp từ hình ảnh (DICDM) sử dụng Extreme Learning Machine (ELM).
    \item \textbf{Ưu điểm:} 
    \begin{itemize}
        \item FADM sử dụng KPCA cho độ chính xác rất cao (lên đến 98.99\%), giúp phân loại rất tốt nhờ trích xuất đặc trưng hình thái cụ thể.
        \item Xử lý giảm thiểu tình trạng overfitting thông qua Cross-Validation.
    \end{itemize}
    \item \textbf{Nhược điểm:} 
    \begin{itemize}
        \item Cần một quy trình tiền xử lý và trích xuất đặc trưng thủ công, phức tạp và tốn thời gian.
        \item Tập dữ liệu thử nghiệm trong bài vẫn ở quy mô tương đối nhỏ so với thực tế.
    \end{itemize}
\end{itemize}

\textbf{Bài báo 2: Empowering Agricultural Insights: RiceLeafBD - A Novel Dataset and Optimal Model Selection for Rice Leaf Disease Diagnosis through Transfer Learning Technique}
\begin{itemize}
    \item \textbf{Mô hình sử dụng:} Phân tích dựa trên các mô hình Học sâu (Deep Learning) và Học chuyển giao (Transfer Learning) bao gồm Light CNN, InceptionNet-V2, MobileNet-V2, và EfficientNet-V2.
    \item \textbf{Ưu điểm:} 
    \begin{itemize}
        \item Mô hình EfficientNet-V2 đạt độ chính xác ấn tượng (91.5\%) trên bộ dữ liệu thực tế (RiceLeafBD) thu thập tại Bangladesh (không phải ảnh lab).
        \item Khả năng ứng dụng thực tiễn cao nhờ việc học các đặc trưng phức tạp từ dữ liệu đa dạng (các điều kiện ánh sáng, nền khác nhau).
    \end{itemize}
    \item \textbf{Nhược điểm:} 
    \begin{itemize}
        \item Các mô hình Transfer Learning (đặc biệt InceptionNet-V2) yêu cầu tài nguyên tính toán khá lớn và có thể khó triển khai trực tiếp trên các thiết bị di động có cấu hình yếu nếu không tiến hành tối ưu hóa.
        \item Cần nhiều thời gian cho việc huấn luyện (training).
    \end{itemize}
\end{itemize}

\textbf{Bài báo 3: Advancing Green AI: Efficient and Accurate Lightweight CNNs for Rice Leaf Disease Identification}
\begin{itemize}
    \item \textbf{Mô hình sử dụng:} Đánh giá các mạng nơ-ron tích chập (CNN) cấu trúc nhẹ và hướng tới "Green AI" như ShuffleNet, MobileNetV2, và EfficientNet-B0.
    \item \textbf{Ưu điểm:} 
    \begin{itemize}
        \item Tập trung vào các mô hình có kích thước nhỏ, tiết kiệm năng lượng và tài nguyên tính toán, rất phù hợp triển khai trên các thiết bị di động (edge devices) hoặc IoT hệ thống theo dõi thực thời.
        \item EfficientNet-B0 đạt hiệu năng rất cao (99.8\%) với kích thước bộ tham số chỉ 5.4 triệu.
    \end{itemize}
    \item \textbf{Nhược điểm:} 
    \begin{itemize}
        \item Rất dễ bị overfitting do thiết lập learning rate và epochs nếu không có cơ chế early stopping hợp lý.
        \item Một số mô hình cực nhẹ như ShuffleNet lại đạt độ chính xác khá thấp (chỉ khoảng 66\%), không đáp ứng được yêu cầu của hệ thống chuẩn đoán thực tế.
    \end{itemize}
\end{itemize}

\begin{figure}[htbp]
    \centering
    % Placeholder for Figure 1. Inster image extracted from references/b3.pdf Page 4
    \framebox[\textwidth]{\parbox[c][5cm][c]{\textwidth}{\centering \textbf{\huge [Chèn Hình Ảnh Từ b3.pdf Trang 4 Vào Đây]}}}
    \caption{Ví dụ về minh họa mô hình nhận diện trong Bài báo 1 (b3.pdf)}
    \label{fig:b3_image}
\end{figure}

\begin{figure}[htbp]
    \centering
    % Placeholder for Figure 2. Inster image extracted from references/b4.pdf Page 9
    \framebox[\textwidth]{\parbox[c][5cm][c]{\textwidth}{\centering \textbf{\huge [Chèn Hình Ảnh Từ b4.pdf Trang 9 Vào Đây]}}}
    \caption{Sơ đồ kiến trúc mô hình Light CNN trích từ Bài báo 2 (b4.pdf)}
    \label{fig:b4_image}
\end{figure}

\subsection{2. Mô hình mà đề tài sử dụng và lý do lựa chọn}

\textbf{Mô hình sử dụng:} 
Đề tài sử dụng mô hình \textbf{prithivMLmods/Rice-Leaf-Disease}, một mô hình tiên tiến thuộc dòng \textbf{Vision Transformer (ViT)}, cụ thể được tinh chỉnh dựa trên kiến trúc gốc vit-base-patch16-224 của Google. Mô hình này được nhúng vào một \textbf{Hệ thống AI đa tác nhân (Multi-Agent System)} bao gồm LLM và Vector Database (ChromaDB/Qdrant) để phối hợp chẩn đoán và giải thích.

\textbf{Lý do sử dụng:}
\begin{itemize}
    \item \textbf{Lợi thế của Vision Transformer (ViT):} Khác với CNN truyền thống (tập trung vào các đặc trưng cục bộ do phép tích chập - convolution), ViT vận hành dựa trên \textit{cơ chế Tự chú ý (Self-Attention)}. Khi chia ảnh thành các ô vuông nhỏ (patches), mạng có thể học được "mối quan hệ tương quan toàn cục" giữa các vùng khác nhau trên lá lúa. Điều này giúp ViT có khả năng nhận diện xuất sắc các loại bệnh có biểu hiện hình thái phức tạp, cấu trúc trải dài trên bề mặt lá thay vì chỉ các đốm riêng lẻ.
    \item \textbf{Giải quyết bài toán "Hộp đen" (Black-box) của Deep Learning:} Mặc dù ViT cung cấp độ chính xác phân loại cao, nhưng người dùng nông nghiệp khó hiểu lý do đưa ra kết luận. Do đó, việc sử dụng hệ thống Multi-Agent có thể bù đắp khiếm khuyết này:
    \begin{itemize}
        \item \textit{Agent Phân tích Hình thái:} Xác minh lại các nhận xét về hình dạng, màu sắc để củng cố độ tin cậy.
        \item \textit{Agent Tra cứu (RAG):} Kết nối Vector DB để đưa ra giải thích bệnh học (tác nhân nấm/vi khuẩn) và kết hợp với LLM biên dịch thành các lời khuyên phòng ngừa sinh học dễ hiểu.
    \end{itemize}
    \item \textbf{Hiệu suất cao trên dữ liệu chuẩn hóa:} Kết hợp với bộ dữ liệu sharmin3 chất lượng, hệ thống đảm bảo cho ra kết luận chẩn đoán tin cậy và trực quan cho chuyên gia nông nghiệp.
\end{itemize}


\newpage
\section{Task 2: Trình bày chi tiết về kiến trúc của mô hình}

\subsection{Tổng quan kiến trúc Framework}

Hệ thống được xây dựng theo mô hình \textbf{Đa tác nhân (Multi-Agent System)} --- một kiến trúc phần mềm trong đó nhiều tác nhân AI (agents) hoạt động độc lập và cộng tác với nhau để giải quyết một bài toán phức tạp hơn khả năng của bất kỳ agent đơn lẻ nào. Thay vì dùng một mạng học sâu "black-box" duy nhất đưa ra kết luận, hệ thống Multi-Agent phân tách bài toán thành các bước chuyên biệt và minh bạch theo từng agent riêng lẻ.

\textbf{Nền tảng công nghệ (Technology Stack):}
\begin{itemize}
    \item \textbf{LangChain / LlamaIndex:} Framework điều phối các agent và quản lý chuỗi gọi công cụ (tool-calling chain).
    \item \textbf{Vision Transformer (ViT):} Mô hình học sâu cốt lõi \textit{prithivMLmods/Rice-Leaf-Disease} từ HuggingFace dùng để nhận diện và phân loại bệnh lúa.
    \item \textbf{LLM (GPT-4o / Gemini):} Mô hình ngôn ngữ lớn đóng vai trò "bộ não" của Agent Điều phối và Agent Tra cứu, giúp diễn giải và tổng hợp thông tin thành ngôn ngữ tự nhiên dễ hiểu.
    \item \textbf{Vector Database (Qdrant / ChromaDB):} Cơ sở dữ liệu vector lưu trữ kiến thức bệnh học được mã hóa (embeddings) từ các tài liệu nông nghiệp, sách giáo khoa bệnh cây trồng, và hướng dẫn phòng trị.
    \item \textbf{FastAPI / Gradio:} Giao diện backend nhận ảnh từ người dùng và trả về kết quả chẩn đoán.
\end{itemize}

\textbf{Cơ chế giao tiếp giữa các Agent:}

Các agent trong hệ thống giao tiếp với nhau thông qua một \textit{Shared Memory / Message Bus} do Agent Điều phối (Manager Agent) kiểm soát. Cụ thể:
\begin{enumerate}
    \item Agent Điều phối nhận yêu cầu từ người dùng và phân phối tác vụ cho các agent con bằng cách gọi đến các \textit{Tool Function} tương ứng.
    \item Mỗi agent con sau khi hoàn thành nhiệm vụ sẽ trả kết quả về dưới dạng một chuỗi văn bản có cấu trúc (structured output).
    \item Agent Điều phối thu thập tất cả các kết quả, đưa chúng vào \textit{context window} của LLM và yêu cầu LLM tổng hợp thành một báo cáo chẩn đoán hoàn chỉnh, súc tích bằng ngôn ngữ tự nhiên.
\end{enumerate}

Cơ chế \textbf{Retrieval-Augmented Generation (RAG)} cho phép Agent Tra cứu không bị giới hạn bởi kiến thức nội tại của LLM. Thay vào đó, mỗi khi cần thông tin về một loại bệnh cụ thể, agent sẽ vector-hóa tên bệnh thành một \textit{query embedding} và thực hiện tìm kiếm ngữ nghĩa (semantic search) trong Vector Database để truy xuất đúng đoạn văn bản bệnh học liên quan nhất. Điều này đảm bảo tính chính xác và cập nhật của thông tin ngay cả khi LLM không "nhớ" đầy đủ dữ liệu chuyên ngành.

\subsection{Quy trình và Vai trò của từng Agent}

\begin{enumerate}[label=\textbf{B\arabic*:}]
    \item \textbf{Agent Tiếp nhận và Điều phối (Manager Agent):}
    Nhận yêu cầu và hình ảnh từ người dùng qua Frontend (Web/App). Kích hoạt và theo dõi luồng công việc của các tác nhân con. Cuối cùng, tổng hợp mọi kết quả từ các Agent khác để trả về báo cáo phân tích trực quan.
    
    \item \textbf{Agent Tiền xử lý dữ liệu:}
    Thực hiện làm sạch, loại bỏ nhiễu, cân bằng sáng/tương phản. Đặc biệt là tiến hành \textit{Tách nền} (loại bỏ tay người, đất, cỏ dại) để các model AI sau chỉ tập trung nhận diện vùng pixels của lá lúa.
    
    \item \textbf{Agent Phân loại bệnh (Tính toán song song - Luồng số 1):}
    Sử dụng mô hình cốt lõi Vision Transformer (prithivMLmods/Rice-Leaf-Disease). Input ảnh 224x224 được chia thành các patches 16x16, chạy qua khối Self-Attention để xuất ra tên bệnh lý (Ví dụ: Đạo ôn, Bạc lá, Đốm nâu...) kèm theo xác suất độ tin cậy (Confidence Score).
    
    \item \textbf{Agent Phân tích Hình thái (Tính toán song song - Luồng số 2):}
    Đóng vai trò cơ chế "Double Check". Model nhỏ hơn sẽ phân tích các \textit{đặc trưng thị giác} như: Vết bệnh màu gì (nâu, xám, vàng...)? Hình dạng thế nào (thoi, vệt dài...)? Có tính đối xứng/viền mờ không? Việc này nhằm so khớp, đảm bảo logic với nhãn bệnh mà Agent Phân loại đã gọi tên.
    
    \item \textbf{Agent Tra cứu Tri thức (Retrieval-Augmented Generation - RAG):}
    Sau khi có kết luận chuẩn xác, tác nhân này sẽ dùng công cụ LLM (GPT/Gemini) chèn Tên bệnh lý dạng Query vào Cơ sở dữ liệu Vector (ChromaDB hoặc Qdrant). Nội dung trả về sẽ là: tác nhân gây hại (nấm, vi rút), điều kiện phát triển lý tưởng và khuyến nghị phương pháp bảo vệ thực vật sinh học hiệu quả.
\end{enumerate}

\subsection{Sơ đồ Kiến trúc Hệ thống Multi-Agent}

\vspace{0.5cm}
\begin{center}
\begin{tikzpicture}[
    node distance=1.5cm and 1cm,
    base/.style={rectangle, draw, rounded corners, text centered, minimum height=1.2cm},
    user/.style={base, fill=gray!10, minimum width=4cm},
    mainbox/.style={base, fill=orange!20, font=\bfseries, minimum width=4cm},
    agent/.style={base, fill=blue!10, minimum width=3.5cm, text width=3.5cm},
    db/.style={cylinder, shape border rotate=90, draw, fill=green!10, minimum height=1cm, text centered},
    arrow/.style={thick,->,>=stealth},
    line/.style={thick,-}
]

% Nodes
\node (user) [user] {Người Dùng (Input Image)};
\node (coord) [mainbox, below=1.5cm of user] {Agent Điều phối};
\node (preproc) [agent, right=2cm of coord] {Agent Tiền xử lý};

% Split nodes
\node (classify) [agent, below=2cm of preproc, xshift=-2.8cm] {Agent Phân loại (ViT)};
\node (morph) [agent, below=2cm of preproc, xshift=2.8cm] {Agent Phân tích Hình thái};

\node (knowledge) [agent, below=2cm of classify] {Agent Tra cứu (RAG)};
\node (vectordb) [db, right=2cm of knowledge] {Vector DB (Qdrant)};
\node (output) [user, below=7cm of coord] {Báo cáo chẩn đoán (Frontend)};

% Connections
\draw [arrow] (user) -- (coord);
\draw [arrow] (coord) -- node[above] {Kích hoạt} (preproc);

\draw [arrow] (preproc) -- ++(0,-1cm) -| node[above, pos=0.25] {Ảnh sạch} (classify);
\draw [arrow] (preproc) -- ++(0,-1cm) -| node[above, pos=0.25] {Ảnh sạch} (morph);

\draw [arrow] (morph) -- node[above] {So khớp chéo} (classify);
\draw [arrow] (classify) -- node[left] {Nhãn bệnh} (knowledge);

\draw [thick,<->,>=stealth] (knowledge) -- node[above] {Query/Return} (vectordb);

% Return paths
\draw [arrow] (knowledge) |- (coord);
\draw [arrow, dashed] (classify) -- (-4.5, -4.5) |- (coord);
\draw [arrow] (coord) -- (output);

\end{tikzpicture}
\end{center}
\vspace{0.5cm}
\textit{Hình 3: Biểu đồ luồng công việc Hệ thống Đa tác nhân (Multi-Agent Workflow)}

\end{document}
